\chapter{Communication Complexity — Public Coin Finite Message}
%%% generated by the blueprint pipeline from TCSlib.CommunicationComplexity.NewmanTheorem.PublicCoinFiniteMessage
\section{Overview}

This module defines public-coin finite-message protocols, in which Alice and Bob each
receive a shared random string $\omega \in \Omega$ alongside their private inputs.
It introduces the notions of approximate satisfaction of a predicate and approximate
computation of a function, and establishes equivalences between the finite-message and
binary public-coin protocol representations, preserving both run behavior and complexity.

\section{Declarations}

\begin{definition}[Public-coin finite-message protocol]
\lean{CommunicationComplexity.PublicCoin.FiniteMessage.Protocol}
\label{CommunicationComplexity.PublicCoin.FiniteMessage.Protocol}
\leanok
\uses{CommunicationComplexity.Deterministic.FiniteMessage.Protocol}
A public-coin finite-message protocol over shared-randomness space $\Omega$, private input
sets $X$ and $Y$, and output type $\alpha$ is a deterministic finite-message protocol in
which Alice's effective input is $\Omega \times X$ and Bob's effective input is
$\Omega \times Y$.
\end{definition}

\begin{definition}[Output node]
\lean{CommunicationComplexity.PublicCoin.FiniteMessage.Protocol.output}
\label{CommunicationComplexity.PublicCoin.FiniteMessage.Protocol.output}
\leanok
\uses{CommunicationComplexity.PublicCoin.FiniteMessage.Protocol}
The terminal (output) node of a public-coin finite-message protocol returning the
constant value $a : \alpha$.
\end{definition}

\begin{definition}[Alice's send step]
\lean{CommunicationComplexity.PublicCoin.FiniteMessage.Protocol.alice}
\label{CommunicationComplexity.PublicCoin.FiniteMessage.Protocol.alice}
\leanok
\uses{CommunicationComplexity.PublicCoin.FiniteMessage.Protocol}
Given a message function $f : X \to \Omega \to \beta$ and a continuation
$P : \beta \to \mathrm{Protocol}$, this constructs the protocol node in which Alice sends
a $\beta$-valued message determined by her private input $x$ and the shared randomness
$\omega$, then continues with $P$.
\end{definition}

\begin{definition}[Bob's send step]
\lean{CommunicationComplexity.PublicCoin.FiniteMessage.Protocol.bob}
\label{CommunicationComplexity.PublicCoin.FiniteMessage.Protocol.bob}
\leanok
\uses{CommunicationComplexity.PublicCoin.FiniteMessage.Protocol}
Given a message function $f : Y \to \Omega \to \beta$ and a continuation
$P : \beta \to \mathrm{Protocol}$, this constructs the protocol node in which Bob sends
a $\beta$-valued message determined by his private input $y$ and the shared randomness
$\omega$, then continues with $P$.
\end{definition}

\begin{definition}[Randomized execution]
\lean{CommunicationComplexity.PublicCoin.FiniteMessage.Protocol.rrun}
\label{CommunicationComplexity.PublicCoin.FiniteMessage.Protocol.rrun}
\leanok
\uses{CommunicationComplexity.Deterministic.FiniteMessage.Protocol.run, CommunicationComplexity.PublicCoin.FiniteMessage.Protocol}
$\mathrm{rrun}(p, x, y, \omega)$ executes the protocol $p$ on private inputs $x$ and $y$
with shared randomness $\omega$, defined as $p.\mathrm{run}(\omega, x)\,(\omega, y)$.
\end{definition}

\begin{theorem}[Randomized execution unfolds to deterministic execution]
\lean{CommunicationComplexity.PublicCoin.FiniteMessage.Protocol.rrun_eq}
\label{CommunicationComplexity.PublicCoin.FiniteMessage.Protocol.rrun_eq}
\leanok
\uses{CommunicationComplexity.Deterministic.FiniteMessage.Protocol.run, CommunicationComplexity.PublicCoin.FiniteMessage.Protocol, CommunicationComplexity.PublicCoin.FiniteMessage.Protocol.rrun}
\difficulty{1}
Let $p$ be a public-coin finite-message protocol over shared-randomness space $\Omega$,
private input sets $X$ and $Y$, and output type $\alpha$, and let $x \in X$, $y \in Y$,
and $\omega \in \Omega$. Then the randomized execution of $p$ on inputs $x$ and $y$ with
shared randomness $\omega$ agrees with the deterministic execution of $p$ on Alice's
effective input $(\omega, x)$ and Bob's effective input $(\omega, y)$: \[
\mathrm{rrun}(p, x, y, \omega) = \mathrm{run}(p, (\omega, x), (\omega, y)). \]
\end{theorem}

\begin{definition}[Approximate satisfaction of a predicate]
\lean{CommunicationComplexity.PublicCoin.FiniteMessage.Protocol.ApproxSatisfies}
\label{CommunicationComplexity.PublicCoin.FiniteMessage.Protocol.ApproxSatisfies}
\leanok
\uses{CommunicationComplexity.PublicCoin.FiniteMessage.Protocol, CommunicationComplexity.PublicCoin.FiniteMessage.Protocol.rrun}
A protocol $p$ $\varepsilon$-satisfies a predicate $Q : X \to Y \to \alpha \to \mathrm{Prop}$
if for every input pair $(x, y)$,
\[
  \Pr_{\omega}\bigl[\neg Q\,x\,y\,(p.\mathrm{rrun}\,x\,y\,\omega)\bigr] \;\le\; \varepsilon.
\]
\end{definition}

\begin{definition}[Approximate computation of a function]
\lean{CommunicationComplexity.PublicCoin.FiniteMessage.Protocol.ApproxComputes}
\label{CommunicationComplexity.PublicCoin.FiniteMessage.Protocol.ApproxComputes}
\leanok
\uses{CommunicationComplexity.PublicCoin.FiniteMessage.Protocol, CommunicationComplexity.PublicCoin.FiniteMessage.Protocol.rrun}
A protocol $p$ $\varepsilon$-computes a function $f : X \to Y \to \alpha$ if for every
input pair $(x, y)$,
\[
  \Pr_{\omega}\bigl[p.\mathrm{rrun}\,x\,y\,\omega \ne f\,x\,y\bigr] \;\le\; \varepsilon.
\]
\end{definition}

\begin{theorem}[Approximate computation as approximate satisfaction of equality]
\lean{CommunicationComplexity.PublicCoin.FiniteMessage.Protocol.ApproxComputes_eq_ApproxSatisfies}
\label{CommunicationComplexity.PublicCoin.FiniteMessage.Protocol.ApproxComputes_eq_ApproxSatisfies}
\leanok
\uses{CommunicationComplexity.PublicCoin.FiniteMessage.Protocol, CommunicationComplexity.PublicCoin.FiniteMessage.Protocol.ApproxComputes, CommunicationComplexity.PublicCoin.FiniteMessage.Protocol.ApproxSatisfies}
\difficulty{1}
Let $\Omega$ be a probability space, and let $p$ be a public-coin finite-message
protocol with shared randomness in $\Omega$, private inputs from $X$ and $Y$, and
outputs in $\alpha$. For every function $f : X \to Y \to \alpha$ and every real number
$\varepsilon$, the assertion that $p$ $\varepsilon$-computes $f$ coincides with the
assertion that $p$ $\varepsilon$-satisfies the predicate $(x, y, a) \mapsto (a =
f\,x\,y)$; that is, the two propositions are equal.
\end{theorem}

\begin{definition}[Conversion to binary public-coin protocol]
\lean{CommunicationComplexity.PublicCoin.FiniteMessage.Protocol.toProtocol}
\label{CommunicationComplexity.PublicCoin.FiniteMessage.Protocol.toProtocol}
\leanok
\uses{CommunicationComplexity.Deterministic.FiniteMessage.Protocol.toProtocol, CommunicationComplexity.PublicCoin.FiniteMessage.Protocol, CommunicationComplexity.PublicCoin.Protocol}
Converts a public-coin finite-message protocol to a binary public-coin protocol by
delegating to \texttt{CommunicationComplexity.Deterministic.FiniteMessage.Protocol.toProtocol}.
\end{definition}

\begin{theorem}[Binary conversion preserves randomized execution]
\lean{CommunicationComplexity.PublicCoin.FiniteMessage.Protocol.toProtocol_rrun}
\label{CommunicationComplexity.PublicCoin.FiniteMessage.Protocol.toProtocol_rrun}
\leanok
\uses{CommunicationComplexity.Deterministic.FiniteMessage.Protocol.toProtocol_run, CommunicationComplexity.PublicCoin.FiniteMessage.Protocol, CommunicationComplexity.PublicCoin.FiniteMessage.Protocol.rrun, CommunicationComplexity.PublicCoin.FiniteMessage.Protocol.toProtocol, CommunicationComplexity.PublicCoin.Protocol.rrun}
\difficulty{1}
Let $p$ be a public-coin finite-message protocol with shared-randomness space $\Omega$,
private input sets $X$ and $Y$, and output type $\alpha$, and let $q$ be the binary
public-coin protocol obtained from $p$ by encoding each finite-alphabet message as a
string of bits. Then for every pair of private inputs $x \in X$ and $y \in Y$ and every
shared random string $\omega \in \Omega$, the two protocols produce the same output when
executed on those inputs with that randomness: the randomized execution of $q$ on $x$,
$y$, $\omega$ equals the randomized execution of $p$ on $x$, $y$, $\omega$.
\end{theorem}

\begin{theorem}[Conversion to a binary public-coin protocol preserves communication complexity]
\lean{CommunicationComplexity.PublicCoin.FiniteMessage.Protocol.toProtocol_complexity}
\label{CommunicationComplexity.PublicCoin.FiniteMessage.Protocol.toProtocol_complexity}
\leanok
\uses{CommunicationComplexity.Deterministic.FiniteMessage.Protocol.complexity, CommunicationComplexity.Deterministic.FiniteMessage.Protocol.toProtocol_complexity, CommunicationComplexity.Deterministic.Protocol.complexity, CommunicationComplexity.PublicCoin.FiniteMessage.Protocol, CommunicationComplexity.PublicCoin.FiniteMessage.Protocol.toProtocol}
\difficulty{1}
Let $\Omega$ be a shared-randomness space and let $p$ be a public-coin finite-message
protocol with Alice's private input in $X$, Bob's in $Y$, and output in $\alpha$; that
is, a deterministic finite-message protocol whose effective input types are $\Omega
\times X$ for Alice and $\Omega \times Y$ for Bob, in which each message is an element
of an arbitrary finite nonempty type $\beta$ costing $\lceil \log_2 \abs{\beta} \rceil$
bits. Then the binary public-coin protocol obtained from $p$ by encoding every
$\beta$-valued message as $\lceil \log_2 \abs{\beta} \rceil$ bits has communication
complexity exactly equal to that of $p$.
\end{theorem}

\begin{definition}[Embedding of binary public-coin protocol]
\lean{CommunicationComplexity.PublicCoin.FiniteMessage.Protocol.ofProtocol}
\label{CommunicationComplexity.PublicCoin.FiniteMessage.Protocol.ofProtocol}
\leanok
\uses{CommunicationComplexity.Deterministic.FiniteMessage.Protocol.ofProtocol, CommunicationComplexity.PublicCoin.FiniteMessage.Protocol, CommunicationComplexity.PublicCoin.Protocol}
Embeds a binary public-coin protocol into a finite-message protocol by delegating to
\texttt{CommunicationComplexity.Deterministic.FiniteMessage.Protocol.ofProtocol}.
\end{definition}

\begin{theorem}[Embedding preserves randomized execution]
\lean{CommunicationComplexity.PublicCoin.FiniteMessage.Protocol.ofProtocol_rrun}
\label{CommunicationComplexity.PublicCoin.FiniteMessage.Protocol.ofProtocol_rrun}
\leanok
\uses{CommunicationComplexity.Deterministic.FiniteMessage.Protocol.ofProtocol_run, CommunicationComplexity.PublicCoin.FiniteMessage.Protocol.ofProtocol, CommunicationComplexity.PublicCoin.FiniteMessage.Protocol.rrun, CommunicationComplexity.PublicCoin.Protocol, CommunicationComplexity.PublicCoin.Protocol.rrun}
\difficulty{1}
Let $\Omega$ be a shared-randomness space, let $X$ and $Y$ be the private input sets of
the two players, and let $\alpha$ be an output type. Let $p$ be a binary public-coin
protocol over these types, that is, a deterministic protocol whose players receive
inputs in $\Omega \times X$ and $\Omega \times Y$. Embedding $p$ into the finite-message
framework—by treating each Boolean message as a message drawn from the two-element
type—and then running the result yields the same output as running $p$ directly: for
every private input $x \in X$, every private input $y \in Y$, and every shared random
string $\omega \in \Omega$, the finite-message embedding of $p$, executed on the paired
inputs $(\omega, x)$ and $(\omega, y)$, returns the same value in $\alpha$ as $p$
executed on those same paired inputs.
\end{theorem}

\begin{theorem}[Complexity invariance of the finite-message embedding]
\lean{CommunicationComplexity.PublicCoin.FiniteMessage.Protocol.ofProtocol_complexity}
\label{CommunicationComplexity.PublicCoin.FiniteMessage.Protocol.ofProtocol_complexity}
\leanok
\uses{CommunicationComplexity.Deterministic.FiniteMessage.Protocol.complexity, CommunicationComplexity.Deterministic.FiniteMessage.Protocol.ofProtocol_complexity, CommunicationComplexity.Deterministic.Protocol.complexity, CommunicationComplexity.PublicCoin.FiniteMessage.Protocol.ofProtocol, CommunicationComplexity.PublicCoin.Protocol}
\difficulty{1}
Let $\Omega$ be a shared-randomness space, let $X$ and $Y$ be the private input sets of
Alice and Bob, let $\alpha$ be an output type, and let $p$ be a binary public-coin
protocol over shared randomness $\Omega$ with these private inputs and this output —
that is, a deterministic protocol whose players branch on Boolean messages, with Alice
reading $\Omega \times X$ and Bob reading $\Omega \times Y$. Embedding $p$ into the
finite-message framework by regarding each Boolean message as an element of the
two-element type $\mathrm{Bool}$ yields a finite-message protocol whose worst-case
communication complexity equals the worst-case communication complexity of $p$.
\end{theorem}

\begin{theorem}[Every public-coin protocol has an equal-complexity finite-message form]
\lean{CommunicationComplexity.PublicCoin.FiniteMessage.Protocol.ofProtocol_equiv}
\label{CommunicationComplexity.PublicCoin.FiniteMessage.Protocol.ofProtocol_equiv}
\leanok
\uses{CommunicationComplexity.Deterministic.FiniteMessage.Protocol.complexity, CommunicationComplexity.Deterministic.Protocol.complexity, CommunicationComplexity.PublicCoin.FiniteMessage.Protocol, CommunicationComplexity.PublicCoin.FiniteMessage.Protocol.ofProtocol, CommunicationComplexity.PublicCoin.FiniteMessage.Protocol.ofProtocol_complexity, CommunicationComplexity.PublicCoin.FiniteMessage.Protocol.ofProtocol_rrun, CommunicationComplexity.PublicCoin.FiniteMessage.Protocol.rrun, CommunicationComplexity.PublicCoin.Protocol, CommunicationComplexity.PublicCoin.Protocol.rrun}
\difficulty{1}
For every public-coin protocol $p$ over shared-randomness space $\Omega$, private input
sets $X$ and $Y$, and output type $\alpha$, there exists a public-coin finite-message
protocol $P$ over the same data that computes exactly the same output under randomized
execution and has the same communication complexity. Precisely, $\mathrm{rrun}(P, x, y,
\omega) = \mathrm{rrun}(p, x, y, \omega)$ for all private inputs $x \in X$, $y \in Y$
and every shared random string $\omega \in \Omega$, and the communication complexity of
$P$ equals that of $p$.
\end{theorem}
